\section{Results}
\label{sec:results}

\subsection{Sealed physical efficacy}

% AUTO-GENERATED by analyze_study2_secondary.py. DO NOT EDIT.
\begin{table*}[t]
\centering
\caption{Headline sealed Study 2 comparison over the same
\studyclaim{StudyTwoDesignCount} designs. Penalized
Fmax is averaged equally over designs and scores incorrect or implementation-
failed samples as zero. Confidence intervals are paired, stratified 95\%
bootstrap intervals versus the SFT baseline.}
\label{tab:study2-primary}
\small
\begin{tabular}{lrrrr}
\toprule
Arm & Training seeds & Penalized Fmax [MHz] & Correct [\%] & $\Delta$ vs. SFT [95\% CI] \\
\midrule
SFT & -- & \studyclaim{StudyTwoOverallSftPenalized} &
\studyclaim{StudyTwoSftCorrectness} & -- \\
Correctness only & \studyclaim{StudyTwoCorrectnessTrainingSeeds} &
\studyclaim{StudyTwoOverallCorrectnessOnlyPenalized} &
\studyclaim{StudyTwoCorrectnessOnlyCorrectness} &
\studyclaim{StudyTwoOverallCorrectnessOnlySftGain}
[\studyclaim{StudyTwoOverallCorrectnessOnlySftCiLow},
 \studyclaim{StudyTwoOverallCorrectnessOnlySftCiHigh}] \\
Original MLP reward & \studyclaim{StudyTwoMlpTrainingSeeds} &
\studyclaim{StudyTwoOverallMlpPenalized} & \studyclaim{StudyTwoMlpCorrectness} &
\studyclaim{StudyTwoOverallMlpSftGain}
[\studyclaim{StudyTwoOverallMlpSftCiLow},
 \studyclaim{StudyTwoOverallMlpSftCiHigh}] \\
Repaired RF reward & \studyclaim{StudyTwoRfTrainingSeeds} &
\studyclaim{StudyTwoOverallRfPenalized} & \studyclaim{StudyTwoRfCorrectness} &
\studyclaim{StudyTwoOverallRfSftGain}
[\studyclaim{StudyTwoOverallRfSftCiLow},
 \studyclaim{StudyTwoOverallRfSftCiHigh}] \\
\bottomrule
\end{tabular}
\end{table*}


Figure~\ref{fig:study2-primary} shows every design and the frozen paired
analysis.  Across \studyclaim{StudyTwoDesignCount} held-out designs, the
repaired-RF arm raises penalized equal-sample \fmax{} from
\studyclaim{StudyTwoOverallSftPenalized} to
\studyclaim{StudyTwoOverallRfPenalized}~MHz.  The paired gain is
\studyclaim{StudyTwoOverallRfSftGain}~MHz with interval
[\studyclaim{StudyTwoOverallRfSftCiLow},
\studyclaim{StudyTwoOverallRfSftCiHigh}]~MHz.

The result is not confined to the easier stratum.  Interpolation improves by
\studyclaim{StudyTwoInterpRfSftGain}~MHz with interval
[\studyclaim{StudyTwoInterpRfSftCiLow},
\studyclaim{StudyTwoInterpRfSftCiHigh}], while extrapolation improves by
\studyclaim{StudyTwoExtrapRfSftGain}~MHz with interval
[\studyclaim{StudyTwoExtrapRfSftCiLow},
\studyclaim{StudyTwoExtrapRfSftCiHigh}]~MHz.  All zero-valued designs remain in
Fig.~\ref{fig:study2-primary}; the aggregate is therefore not obtained by
conditioning on successful hardware.

\begin{figure*}[t]
\centering
\includegraphics[width=\textwidth]{fig02_study2_primary}
\caption{Sealed primary efficacy.  Panel (a) retains every design and joins its
SFT and repaired-RF penalized equal-sample post-route \fmax.  Panel (b) reports
the frozen paired effects and the separately displayed training-seed effects.
The overall gain is \studyclaim{StudyTwoOverallRfSftGain}~MHz
[\studyclaim{StudyTwoOverallRfSftCiLow},
\studyclaim{StudyTwoOverallRfSftCiHigh}]~MHz.}
\label{fig:study2-primary}
\end{figure*}

Both independently trained RF policies improve: the seed-specific overall
effects are \studyclaim{StudyTwoRfSeedOneOverallGain} and
\studyclaim{StudyTwoRfSeedTwoOverallGain}~MHz.  Their positive directions also
hold separately under interpolation and extrapolation
(Table~\ref{tab:study2-rf-seeds}).  Correctness changes from
\studyclaim{StudyTwoSftCorrectness}\% for SFT to
\studyclaim{StudyTwoRfCorrectness}\% for RF.  This difference is not hidden:
because every incorrect draw scores zero, it is already paid for in the
primary effect.

% AUTO-GENERATED by analyze_study2_secondary.py. DO NOT EDIT.
\begin{table*}[t]
\centering
\caption{Independent repaired-RF training-seed effects on the sealed endpoint.
Seed rows are separately reported point estimates; the seed-averaged row uses
the frozen paired, stratified 95\% bootstrap intervals.}
\label{tab:study2-rf-seeds}
\small
\begin{tabular}{lrrr}
\toprule
Estimate & Interpolation & Extrapolation & Overall \\
\midrule
RF seed 1 & \studyclaim{StudyTwoRfSeedOneInterpGain} &
\studyclaim{StudyTwoRfSeedOneExtrapGain} &
\studyclaim{StudyTwoRfSeedOneOverallGain} \\
RF seed 2 & \studyclaim{StudyTwoRfSeedTwoInterpGain} &
\studyclaim{StudyTwoRfSeedTwoExtrapGain} &
\studyclaim{StudyTwoRfSeedTwoOverallGain} \\
Seed average [95\% CI] &
\studyclaim{StudyTwoInterpRfSftGain}
[\studyclaim{StudyTwoInterpRfSftCiLow}, \studyclaim{StudyTwoInterpRfSftCiHigh}] &
\studyclaim{StudyTwoExtrapRfSftGain}
[\studyclaim{StudyTwoExtrapRfSftCiLow}, \studyclaim{StudyTwoExtrapRfSftCiHigh}] &
\studyclaim{StudyTwoOverallRfSftGain}
[\studyclaim{StudyTwoOverallRfSftCiLow}, \studyclaim{StudyTwoOverallRfSftCiHigh}] \\
\bottomrule
\end{tabular}
\end{table*}


% AUTO-GENERATED by analyze_study2_secondary.py. DO NOT EDIT.
\begin{table*}[t]
\centering
\caption{Post-primary descriptive family breakdown of the sealed
failure-penalized endpoint. I/E gives interpolation/extrapolation design counts.
Rows are descriptive point estimates; the small and unequal family cells do not
support stable family-level confidence intervals.}
\label{tab:study2-family}
\small
\begin{tabular}{lrrrrr}
\toprule
Family & Designs & I/E & SFT [MHz] & RF [MHz] & RF--SFT [MHz] \\
\midrule
FIR & \studyclaim{StudyTwoFamilyFirCount} & \studyclaim{StudyTwoFamilyFirInterpCount}/\studyclaim{StudyTwoFamilyFirExtrapCount} & \studyclaim{StudyTwoFamilyFirSft} & \studyclaim{StudyTwoFamilyFirRf} & \studyclaim{StudyTwoFamilyFirGain} \\
Reverse FIR & \studyclaim{StudyTwoFamilyReverseFirCount} & \studyclaim{StudyTwoFamilyReverseFirInterpCount}/\studyclaim{StudyTwoFamilyReverseFirExtrapCount} & \studyclaim{StudyTwoFamilyReverseFirSft} & \studyclaim{StudyTwoFamilyReverseFirRf} & \studyclaim{StudyTwoFamilyReverseFirGain} \\
Polynomial & \studyclaim{StudyTwoFamilyPolynomialCount} & \studyclaim{StudyTwoFamilyPolynomialInterpCount}/\studyclaim{StudyTwoFamilyPolynomialExtrapCount} & \studyclaim{StudyTwoFamilyPolynomialSft} & \studyclaim{StudyTwoFamilyPolynomialRf} & \studyclaim{StudyTwoFamilyPolynomialGain} \\
Recursive filter & \studyclaim{StudyTwoFamilyRecursiveCount} & \studyclaim{StudyTwoFamilyRecursiveInterpCount}/\studyclaim{StudyTwoFamilyRecursiveExtrapCount} & \studyclaim{StudyTwoFamilyRecursiveSft} & \studyclaim{StudyTwoFamilyRecursiveRf} & \studyclaim{StudyTwoFamilyRecursiveGain} \\
Median filter & \studyclaim{StudyTwoFamilyMedianCount} & \studyclaim{StudyTwoFamilyMedianInterpCount}/\studyclaim{StudyTwoFamilyMedianExtrapCount} & \studyclaim{StudyTwoFamilyMedianSft} & \studyclaim{StudyTwoFamilyMedianRf} & \studyclaim{StudyTwoFamilyMedianGain} \\
\bottomrule
\end{tabular}
\end{table*}


The post-primary family breakdown prevents the aggregate from becoming a
universal family claim.  FIR, reverse FIR, polynomial, and recursive-filter
point estimates are positive, whereas median is negative and is represented
only by extrapolation designs.  These rows are descriptive because the family
cells are small and unequal.  The primary direction is nevertheless insensitive
to two post-primary checks: the unstratified paired percentile
bootstrap gives
[\studyclaim{StudyTwoSensitivityUnstratifiedCiLow},
\studyclaim{StudyTwoSensitivityUnstratifiedCiHigh}]~MHz, and the mean effect
across leave-one-family-out omissions ranges from
\studyclaim{StudyTwoSensitivityLeaveFamilyOutMin} to
\studyclaim{StudyTwoSensitivityLeaveFamilyOutMax}~MHz.  Both checks are post-
primary sensitivities, not replacements for the frozen interval.

\subsection{Reward-package controls and endpoint decomposition}

% AUTO-GENERATED by paper/make_revision_evidence.py.
\begin{table}[t]
\centering
\caption{Post-primary decomposition on the complete sealed draw set. Conditional frequency pools correct draws, including zero-valued implementation failures.}
\label{tab:decomposition}
\small
\begin{tabular}{@{}lrrr@{}}
\toprule
Policy & $q$ (\%) & $\mu$ (MHz) & $q\mu$ (MHz) \\
\midrule
SFT & \revisionclaim{SftCorrectness} & \revisionclaim{SftConditional} & \revisionclaim{SftPenalized} \\
RF & \revisionclaim{RfCorrectness} & \revisionclaim{RfConditional} & \revisionclaim{RfPenalized} \\
MLP & \revisionclaim{MlpCorrectness} & \revisionclaim{MlpConditional} & \revisionclaim{MlpPenalized} \\
Correctness only & \revisionclaim{CorrectnessCorrectness} & \revisionclaim{CorrectnessConditional} & \revisionclaim{CorrectnessPenalized} \\
\bottomrule
\end{tabular}

\end{table}


Table~\ref{tab:decomposition} separates functional yield from the physical
quality of oracle-passing draws.  With equal budgets across designs, pooled
$q\mu$ reproduces the primary mean exactly; this pooled $\mu$ is not an
unweighted mean of per-design conditional frequencies.  RF increases pooled
conditional frequency from \revisionclaim{SftConditional} to
\revisionclaim{RfConditional}~MHz while correctness decreases.  In contrast,
the correctness-only policy has a conditional frequency of
\revisionclaim{CorrectnessConditional}~MHz.  The observed RF gain is thus a
shift in the physical quality of passing draws, not an increase in their count.

Figure~\ref{fig:study2-controls} separates physical utility from functional
yield.  Correctness-only training reaches
\studyclaim{StudyTwoCorrectnessOnlyCorrectness}\% correctness, above both SFT
and RF, but attains only
\studyclaim{StudyTwoOverallCorrectnessOnlyPenalized}~MHz penalized \fmax.  Its
paired SFT-relative effect is
\studyclaim{StudyTwoOverallCorrectnessOnlySftGain}~MHz with interval
[\studyclaim{StudyTwoOverallCorrectnessOnlySftCiLow},
\studyclaim{StudyTwoOverallCorrectnessOnlySftCiHigh}]~MHz.  Increasing validity
alone therefore does not explain the hardware result in this realized seed.
Its interval resamples designs conditional on that checkpoint; it does not
establish a seed-robust ranking of correctness-only training.

The original MLP reward does improve over SFT, reaching
\studyclaim{StudyTwoOverallMlpPenalized}~MHz with a paired gain of
\studyclaim{StudyTwoOverallMlpSftGain}~MHz.  The repaired RF nonetheless exceeds
it directly by \studyclaim{StudyTwoOverallRfMlpGain}~MHz with interval
[\studyclaim{StudyTwoOverallRfMlpCiLow},
\studyclaim{StudyTwoOverallRfMlpCiHigh}]~MHz.  The RF-minus-MLP intervals remain
above zero under interpolation and extrapolation as well.  Because the arms
share their SFT initialization, tasks, update budget, oracle, and evaluation,
this comparison favors the RF reward package for these checkpoints.  It does
not separate canonicalization, structural features, and predictor architecture,
nor does it compare alternative policy-optimization algorithms.

\begin{figure*}[t]
\centering
\includegraphics[width=\textwidth]{fig03_study2_controls}
\caption{Matched controls.  Panel (a) plots correctness against the
failure-penalized endpoint; the correctness-only arm is the key counterexample.
Panel (b) gives paired contrasts.  Repaired RF exceeds the original MLP reward
by \studyclaim{StudyTwoOverallRfMlpGain}~MHz
[\studyclaim{StudyTwoOverallRfMlpCiLow},
\studyclaim{StudyTwoOverallRfMlpCiHigh}]~MHz.}
\label{fig:study2-controls}
\end{figure*}

\subsection{Conditional physical-cost trade-off}

% AUTO-GENERATED by analyze_study2_secondary.py. DO NOT EDIT.
\begin{table*}[t]
\centering
\caption{Conditional physical-cost diagnostics on the
\studyclaim{StudyTwoPpaCommonSupport}-design three-arm common support. Values
are multiplicity-weighted within design and equally weighted across designs.
Power is Vivado's vectorless estimate, not board power.}
\label{tab:study2-ppa}
\small
\begin{tabular}{lrrrr}
\toprule
Metric & SFT & MLP & RF & RF--SFT difference [CI] \\
\midrule
LUTs & \studyclaim{StudyTwoPpaSftLut} & \studyclaim{StudyTwoPpaMlpLut} & \studyclaim{StudyTwoPpaRfLut} & \studyclaim{StudyTwoPpaRfSftLutDiff} [\studyclaim{StudyTwoPpaRfSftLutCiLow}, \studyclaim{StudyTwoPpaRfSftLutCiHigh}] \\
Flip-flops & \studyclaim{StudyTwoPpaSftFf} & \studyclaim{StudyTwoPpaMlpFf} & \studyclaim{StudyTwoPpaRfFf} & \studyclaim{StudyTwoPpaRfSftFfDiff} [\studyclaim{StudyTwoPpaRfSftFfCiLow}, \studyclaim{StudyTwoPpaRfSftFfCiHigh}] \\
DSP blocks & \studyclaim{StudyTwoPpaSftDsp} & \studyclaim{StudyTwoPpaMlpDsp} & \studyclaim{StudyTwoPpaRfDsp} & \studyclaim{StudyTwoPpaRfSftDspDiff} [\studyclaim{StudyTwoPpaRfSftDspCiLow}, \studyclaim{StudyTwoPpaRfSftDspCiHigh}] \\
BRAM blocks & \studyclaim{StudyTwoPpaSftBram} & \studyclaim{StudyTwoPpaMlpBram} & \studyclaim{StudyTwoPpaRfBram} & \studyclaim{StudyTwoPpaRfSftBramDiff} [\studyclaim{StudyTwoPpaRfSftBramCiLow}, \studyclaim{StudyTwoPpaRfSftBramCiHigh}] \\
Vectorless power [W] & \studyclaim{StudyTwoPpaSftPower} & \studyclaim{StudyTwoPpaMlpPower} & \studyclaim{StudyTwoPpaRfPower} & \studyclaim{StudyTwoPpaRfSftPowerDiff} [\studyclaim{StudyTwoPpaRfSftPowerCiLow}, \studyclaim{StudyTwoPpaRfSftPowerCiHigh}] \\
\bottomrule
\end{tabular}
\end{table*}


The speed improvement is not a free PPA gain.  On the
\studyclaim{StudyTwoPpaCommonSupport}-design common support, RF uses
\studyclaim{StudyTwoPpaRfLut} rather than
\studyclaim{StudyTwoPpaSftLut} LUTs and
\studyclaim{StudyTwoPpaRfFf} rather than
\studyclaim{StudyTwoPpaSftFf} flip-flops.  Conversely, DSP use falls from
\studyclaim{StudyTwoPpaSftDsp} to
\studyclaim{StudyTwoPpaRfDsp} blocks, while BRAM remains unchanged.
Vectorless power rises from \studyclaim{StudyTwoPpaSftPower} to
\studyclaim{StudyTwoPpaRfPower}~W.  The paired intervals in
Table~\ref{tab:study2-ppa} show that the LUT, register, DSP, and vectorless-power
directions are stable on common support.  These conditional quantities answer
what a compiled correct implementation costs; they do not replace the
failure-penalized policy endpoint.

% AUTO-GENERATED by analyze_study2_secondary.py. DO NOT EDIT.
\begin{table*}[t]
\centering
\caption{Illustrative source-level contrast for
\studyclaim{StudyTwoCaseDesign}. The entries are the dominant SFT candidate and
dominant candidate from RF seed one for that design. Selection occurred after
manuscript assessment; this example explains a possible CAD mechanism but is not a
population or causal estimate.}
\label{tab:study2-rtl-case}
\small
\begin{tabularx}{\textwidth}{lXrrrrr}
\toprule
Arm & Observed RTL organization & Sample mass [\%] & Fmax [MHz] & LUT & FF & DSP \\
\midrule
SFT & Delayed samples and product registers followed by a parallel final sum &
\studyclaim{StudyTwoCaseSftMass} & \studyclaim{StudyTwoCaseSftFmax} &
\studyclaim{StudyTwoCaseSftLut} & \studyclaim{StudyTwoCaseSftFf} &
\studyclaim{StudyTwoCaseSftDsp} \\
RF seed one & Registered transposed multiply--accumulate chain &
\studyclaim{StudyTwoCaseRfMass} & \studyclaim{StudyTwoCaseRfFmax} &
\studyclaim{StudyTwoCaseRfLut} & \studyclaim{StudyTwoCaseRfFf} &
\studyclaim{StudyTwoCaseRfDsp} \\
\bottomrule
\end{tabularx}
\end{table*}


Table~\ref{tab:study2-rtl-case} supplies a concrete CAD-level example of that
aggregate trade-off.  For \studyclaim{StudyTwoCaseDesign}, the dominant SFT
candidate stores delayed samples and products before a wide parallel sum.  The
dominant candidate from RF seed one instead emits a registered transposed
multiply--accumulate chain.  Vivado maps the latter to substantially more LUTs
and registers, no DSP blocks, and a much shorter timing path.  Because this
case was selected post-primary for interpretability, it illustrates
a physically meaningful RTL transformation but does not establish that every
family or every gain follows the same transformation.

\subsection{Setup-closure sensitivity}

% AUTO-GENERATED by paper/make_revision_evidence.py.
\begin{table}[t]
\centering
\caption{Descriptive setup-closure sensitivity (MHz). Original: frozen WNS-derived value. Setup: conservative passing endpoint under the pilot flow.}
\label{tab:setup-closure}
\small
\begin{tabular}{@{}lrrrr@{}}
\toprule
Design & SFT orig. & SFT setup & RF orig. & RF setup \\
\midrule
fir42\_v6\_8b & \revisionclaim{ClosureASftOriginal} & \revisionclaim{ClosureASftSetup} & \revisionclaim{ClosureARfOriginal} & \revisionclaim{ClosureARfSetup} \\
firr20\_v6 & \revisionclaim{ClosureBSftOriginal} & \revisionclaim{ClosureBSftSetup} & \revisionclaim{ClosureBRfOriginal} & \revisionclaim{ClosureBRfSetup} \\
iir18\_v7 & \revisionclaim{ClosureCSftOriginal} & \revisionclaim{ClosureCSftSetup} & \revisionclaim{ClosureCRfOriginal} & \revisionclaim{ClosureCRfSetup} \\
med21 & \revisionclaim{ClosureDSftOriginal} & \revisionclaim{ClosureDSftSetup} & \revisionclaim{ClosureDRfOriginal} & \revisionclaim{ClosureDRfSetup} \\
poly6\_v13\_8b & \revisionclaim{ClosureESftOriginal} & \revisionclaim{ClosureESftSetup} & \revisionclaim{ClosureERfOriginal} & \revisionclaim{ClosureERfSetup} \\
\bottomrule
\end{tabular}

\end{table}


The selected-candidate pilot yields Spearman correlation
\revisionclaim{ClosureRho} between original timing-derived and setup-passing
frequencies, with median symmetric percentage error
\revisionclaim{ClosureMedianSmape}\%.  None of the
\revisionclaim{ClosureNonTiedPairs} non-tied pairs reverses direction
(Table~\ref{tab:setup-closure}).  However, the largest candidate error is
\revisionclaim{ClosureWorstSmape}\%: the reverse-FIR RF candidate changes from
\revisionclaim{ClosureBRfOriginal} to
\revisionclaim{ClosureBRfSetup}~MHz.  The median error therefore must not be
read as a uniform accuracy guarantee.  The check supports relative ordering
on this subset, while absolute frequency remains sensitive to implementation
conditions.  Setup passing alone is not full timing sign-off.

\subsection{Descriptive PYNQ-Z2 evidence}

Figure~\ref{fig:board-evidence} reports all paired board outcomes under the
architecture in Fig.~\ref{fig:board-architecture}.  Across
\boardclaim{BoardRunCount} complete sweeps, the maximum
within-entry boundary spread is \boardclaim{BoardMaximumSpread}~MHz.  The echo
canary reaches \boardclaim{BoardCanaryFmax}~MHz and remains at least
\boardclaim{BoardMinimumCanaryMargin}~MHz above every DUT boundary.

\begin{figure}[!t]
\centering
\includegraphics[width=\columnwidth]{fig04b_board_outcomes}
\caption{Measured PYNQ-Z2 clock boundaries under the amended protocol.
All pairs are retained: \boardclaim{BoardRfWins} RF wins,
\boardclaim{BoardTies} ties, and \boardclaim{BoardRfLosses} loss.
Annotations are signed RF-minus-SFT differences in MHz; the dashed line is
the echo-canary boundary. No population interval is inferred from this subset.}
\label{fig:board-evidence}
\end{figure}

% AUTO-GENERATED by analyze_study2_board.py. DO NOT EDIT.
\begin{table}[t]
\centering
\caption{Descriptive PYNQ-Z2 clock-sweep confirmation on
\boardclaim{BoardPairCount} paired designs. Values are medians over
\boardclaim{BoardRunCount} repeated sweeps. The subset required a disclosed
correctness-feasibility amendment and is not independent confirmatory evidence.}
\label{tab:study2-board}
\small
\begin{tabular}{llrrr}
\toprule
Family & Design & SFT & RF & RF--SFT [MHz] \\
\midrule
FIR & fir34\_v1\_8b & \boardclaim{BoardFirSft} & \boardclaim{BoardFirRf} & \boardclaim{BoardFirDifference} \\
FIRR & firr20\_v6 & \boardclaim{BoardFirrSft} & \boardclaim{BoardFirrRf} & \boardclaim{BoardFirrDifference} \\
POLY & poly16\_v13\_8b & \boardclaim{BoardPolySft} & \boardclaim{BoardPolyRf} & \boardclaim{BoardPolyDifference} \\
IIR & iir18\_v7 & \boardclaim{BoardIirSft} & \boardclaim{BoardIirRf} & \boardclaim{BoardIirDifference} \\
MED & med19 & \boardclaim{BoardMedSft} & \boardclaim{BoardMedRf} & \boardclaim{BoardMedDifference} \\
\bottomrule
\end{tabular}
\end{table}


The descriptive mean changes from \boardclaim{BoardSftMean} to
\boardclaim{BoardRfMean}~MHz, a \boardclaim{BoardMeanDifference}-MHz paired
difference.  That mean is driven by two large feed-forward gains.  Polynomial
and IIR tie, while the median pair changes by
\boardclaim{BoardMedDifference}~MHz.  The silicon evidence therefore supports
the main direction heterogeneously; it does not establish a universal family
effect or independently confirm the population estimate.

\subsection{Breadth and backbone transfer}

The earlier, separate \claim{DesignCount}-design evaluation supplies broader
five-family support.  Its penalized endpoint rises from
\claim{OverallSftPenalized} to \claim{OverallGrpoPenalized}~MHz, with
\claim{OverallImproved} of \claim{DesignCount} designs improving.  Feed-forward
families show the most consistent gains; IIR is mixed and median remains slow
in absolute terms.  Because its uncertainty analysis is historical and
post-hoc, this result supports breadth but is not pooled with the sealed study.

The full training recipe also transfers to Qwen-Coder on the earlier
\claim{QwenDesignCount}-design, \claim{QwenFamilyCount}-family universe.
Penalized \fmax{} rises from \claim{QwenOverallSftPenalized} to
\claim{QwenOverallGrpoPenalized}~MHz, a paired gain of
\claim{QwenOverallGain}~MHz with interval
[\claim{QwenOverallCiLow}, \claim{QwenOverallCiHigh}]~MHz; all
\claim{QwenDesignCount} designs improve.  Correctness also rises from
\claim{QwenOverallSftCorrectness}\% to
\claim{QwenOverallGrpoCorrectness}\%.  This reduces a model-specific
explanation, but its narrower family universe prevents direct pooling.

\subsection{Probability reallocation and sampling efficiency}
\label{sec:mechanism}

Repeating the support diagnostic on the sealed samples refines the mechanism
claim.  Among \studyclaim{StudyTwoSealedMechanismCommonCount} designs with
correct sampled strings under both policies, SFT's measured support contains a
candidate at least as fast as the RF conditional mean on
\studyclaim{StudyTwoSealedMechanismSupportCount}.  Mean inverse-Simpson
effective correct-string support contracts from
\studyclaim{StudyTwoSealedMechanismSftEffective} to
\studyclaim{StudyTwoSealedMechanismRfEffective}.  Thus reallocation explains
most common-support cases, but the sealed exceptions prevent the stronger claim
that every optimized implementation was already observed in the finite SFT
sample.  Nor do unseen SFT samples permit a claim of open-ended invention.

On the sealed distribution, one RF draw lies between the empirical SFT
perfect-selector expectations for \revisionclaim{SelectorLowerN} draws
(\revisionclaim{SelectorLowerFmax}~MHz) and
\revisionclaim{SelectorUpperN} draws
(\revisionclaim{SelectorUpperFmax}~MHz).  At the upper integer crossing, the
SFT workload is \revisionclaim{SelectorUpperN} generations and
\revisionclaim{SelectorStreams} oracle simulation streams per design, with
\revisionclaim{SelectorRawImplementations} expected correct-draw implementation
attempts without caching or \revisionclaim{SelectorCachedImplementations}
distinct implementations with content caching.  These are empirical expected
workloads, not measured runtime or speedup.  They exclude policy-training cost
and assume perfect post-route selection on observed support; they do not imply
dominance over an unrestricted SFT search.

\subsection{Contextual bounds}

The same-task API control shows why conditional speed is insufficient.  A
timing-focused prompt raises conditional \fmax{} from
\claim{FrontierPlainConditional} to
\claim{FrontierFastConditional}~MHz while correctness changes from
\claim{FrontierPlainCorrectness}\% to
\claim{FrontierFastCorrectness}\%; its penalized endpoint is
\claim{FrontierFastPenalized}~MHz.  This is one API configuration with a smaller
budget and is not a general frontier-model ranking.

For selected implementation quality, the fastest measured optimized RTL
candidate reaches \claim{HlsExpertGeomean}$\times$ the geometric-mean
throughput of engineer-pragma HLS over \claim{HlsDesignCount} comparable
designs.  It reaches \claim{HlsNaiveGeomean}$\times$ no-pragma HLS.  Because HLS
starts from engineer-written C++ and the RTL value is selected rather than
policy-average, this comparison is context, not the causal baseline.

\subsection{Matched general-RTL specialization cost}

% AUTO-GENERATED by paper/make_revision_evidence.py.
\begin{table}[t]
\centering
\caption{Matched post-primary general-RTL evaluation. Each checkpoint uses the same benchmark, generation settings, and draw budget.}
\label{tab:matched-verilog}
\small
\begin{tabular}{@{}lr@{}}
\toprule
Checkpoint & Pass-at-one (\%) \\
\midrule
SFT & \revisionclaim{VerilogSft} \\
RF seed one & \revisionclaim{VerilogRfOne} \\
RF seed two & \revisionclaim{VerilogRfTwo} \\
\bottomrule
\end{tabular}

\end{table}


The matched VerilogEval extension compares the actual sealed checkpoints on
\revisionclaim{VerilogProblems} official problems with
\revisionclaim{VerilogSamples} draws per checkpoint and problem.  Pass-at-one
is \revisionclaim{VerilogSft}\% for SFT,
\revisionclaim{VerilogRfOne}\% for RF seed one, and
\revisionclaim{VerilogRfTwo}\% for RF seed two
(Table~\ref{tab:matched-verilog}).  The RF-seed-mean paired change is
\revisionclaim{VerilogDelta} percentage points with interval
[\revisionclaim{VerilogCiLow}, \revisionclaim{VerilogCiHigh}].  This indicates
a small additional general-RTL regression, not ``no further degradation.''
The interval resamples problems conditional on the evaluated checkpoints and
does not estimate policy-training seed variance.

The earlier base-model comparison showed pass-at-one decreasing from
\claim{VerilogBasePassone}\% to \claim{VerilogSftPassone}\% after domain SFT.
That experiment uses different generation settings; its absolute values are
not comparable with the matched RF extension and are not pooled.  Together,
the two experiments identify specialization costs at different training stages,
without establishing performance on unseen industrial RTL.

\subsection{Bounded reward-reliability diagnostic}
\label{sec:trajectory}

Figure~\ref{fig:reward-reliability} contrasts the older exploited reward with
the repaired RF trajectory.  In the historical run, early penalized physical
performance rises from \claim{TrajectoryStartPenalized} to
\claim{TrajectoryEarlyPenalized}~MHz.  Late in training, however, the proxy
rises from \claim{TrajectoryLateProxy} to
\claim{TrajectoryFinalProxy}~MHz while correctness falls from
\claim{TrajectoryLateCorrectness}\% to
\claim{TrajectoryFinalCorrectness}\%.  Penalized physical performance
therefore falls from \claim{TrajectoryLatePenalized} to
\claim{TrajectoryFinalPenalized}~MHz.

\begin{figure*}[t]
\centering
\includegraphics[width=\textwidth]{fig06_reward_reliability}
\caption{Bounded reward-reliability evidence.  Panel (a) shows a historical
proxy rising while failure-penalized Vivado utility falls.  Panel (b) shows the
repaired RF midpoint-to-endpoint trajectory: interpolation changes from
\studyclaim{StudyTwoInterpMidFmax} to
\studyclaim{StudyTwoInterpEndFmax}~MHz and extrapolation from
\studyclaim{StudyTwoExtrapMidFmax} to
\studyclaim{StudyTwoExtrapEndFmax}~MHz.}
\label{fig:reward-reliability}
\end{figure*}

A lexical audit identifies one plausible channel in the older reward: a
line-anchored assignment feature increases while its layout-canonical
counterpart does not, and diagnostic rescoring removes
\claim{ReflowRemoved}\% of the raw late reward jump.  The canonical scores are
clamp-saturated and there are \claim{ReflowTokenPairCount} token-identical
cross-checkpoint pairs, so this is bounded attribution rather than a randomized
causal test.  In the repaired run, interpolation remains essentially stable
from midpoint to endpoint and extrapolation improves.  The positive repair
claim comes from the sealed efficacy study; the historical trajectory only
explains why the reporting boundary is necessary.
